\section{Statistical and Machine-Learning Background}
\label{app:theory}
This appendix collects the definitions and results that the implementations rest on and that the main text only uses. The frequentist quantities (Fisher information, Hesse and Minos uncertainties, Neyman-Pearson test, Wilks' theorem, significance) and the Savage-Dickey ratio are in Appendix~\ref{sec:appendix_freq}.

\subsection{Extended Likelihood}
If the total number of candidates $N$ is Poisson-distributed with mean $\nu=N_s+N_b$ and each candidate is drawn independently from the normalized density $f(m,q)=[N_s S+N_b B]/\nu$, the likelihood of the sample is $\frac{\nu^N e^{-\nu}}{N!}\prod_{i}f(m_i,q_i)=\frac{e^{-\nu}}{N!}\prod_i[N_s S(m_i)+N_b B(m_i)]$. The factor $1/N!$ does not depend on the parameters and is dropped, which gives the extended likelihood of Eq.~\eqref{eq:likelihood}. Unlike the likelihood with a fixed total, it constrains the yields $N_s$ and $N_b$ through $N$ itself, and the asymmetry $A$ is estimated together with them.

\subsection{Sampling the Posterior}
The Metropolis-Hastings algorithm builds a Markov chain whose stationary distribution is the target $\pi(\theta)$: from the state $x$ a point $y$ is proposed from $q(y\mid x)$ and accepted with probability $\min\!\big[1,\pi(y)\,q(x\mid y)/(\pi(x)\,q(y\mid x))\big]$, which satisfies detailed balance. Its efficiency depends on how the proposal is matched to the (here correlated) posterior. The ensemble sampler used for the nominal fit~\cite{goodman_weare,emcee} evolves $K$ walkers; to update the walker $X_k$ it picks another walker $X_j$ and proposes $Y=X_j+Z\,(X_k-X_j)$, with $Z\in[1/a,a]$ drawn from $g(z)\propto z^{-1/2}$ ($a=2$), and accepts it with probability $\min\!\big[1,Z^{d-1}\pi(Y)/\pi(X_k)\big]$ ($d$ the number of parameters). The move is invariant under affine transformations of the parameters, so a strongly correlated, elongated posterior is as easy to sample as a spherical one, and it needs no hand-tuned proposal covariance.

Successive states of a chain are correlated. The integrated autocorrelation time of a parameter is $\tau=1+2\sum_{t\geq1}\rho(t)$, with $\rho(t)$ the autocorrelation of the chain at lag $t$, and a chain of $N$ steps contains $N_{\mathrm{eff}}=N/\tau$ effectively independent samples, so that the Monte Carlo error of a posterior mean is $\sigma/\sqrt{N_{\mathrm{eff}}}$. The estimate of $\tau$ is itself reliable only for $N\gtrsim50\,\tau$~\cite{emcee}.

\subsection{Posterior Summaries and Model Comparison}
The credible intervals quoted are the equal-tailed ones (the $16\%$--$84\%$ quantiles for $68\%$); the highest-posterior-density (HPD) interval is the shortest interval with the same probability and coincides with it for a symmetric posterior. A Bayes factor $BF_{ij}=P(D\mid H_i)/P(D\mid H_j)$ is interpreted with the Jeffreys scale~\cite{jeffreys,kass_raftery}: values between $1$ and $3$ are barely worth mentioning, $3$--$10$ substantial, $10$--$30$ strong, $30$--$100$ very strong and above $100$ decisive. The evidence automatically penalizes unnecessary freedom (Ockham's razor): for a Gaussian posterior much narrower than a uniform prior of width $2W$ the Savage-Dickey ratio of $A=0$ is proportional to $W$, so widening the prior of $A$ favors the null hypothesis by construction, which is why the dependence on $W$ is reported.

The posterior predictive check compares the data $y$ with replicas $y^{\mathrm{rep}}$ simulated from the model at parameters drawn from the posterior, through a discrepancy statistic $T$: the posterior predictive $p$-value is $p_B=P\big(T(y^{\mathrm{rep}},\theta)\geq T(y,\theta)\mid y\big)$~\cite{gelman_ppc}. Values close to $0$ or $1$ signal a feature of the data that the model does not reproduce; a statistic that the fit adjusts directly (for example the total asymmetry, when $A$ is a free parameter) is reproduced by construction and tests nothing.

\subsection{Regularization, Bias-Variance and Loss Functions}
For a classifier with output probability $p(x)$ the binary cross-entropy $-\sum_i[y_i\ln p_i+(1-y_i)\ln(1-p_i)]$ is the negative log-likelihood of a Bernoulli distribution, the same principle (minimization of $-\ln L$) as in the mass fits; the Brier score $\frac1N\sum_i(p_i-y_i)^2$~\cite{brier} is the mean squared error of the predicted probability. The logistic regression is regularized by adding a penalty to the negative log-likelihood, $\alpha\sum_j\beta_j^2$ ($L2$, ridge: shrinks all the coefficients) or $\alpha\sum_j|\beta_j|$ ($L1$, lasso~\cite{tibshirani_lasso}: sets some of them exactly to zero), with $\alpha$ chosen by cross-validation (in \texttt{scikit-learn}, $C=1/\alpha$ multiplies the likelihood term). For squared error the expected test error of an estimator $\hat f$ decomposes as $E[(y-\hat f)^2]=\mathrm{Bias}^2+\mathrm{Var}+\sigma^2_{\mathrm{irr}}$: flexible models lower the bias and raise the variance, and the optimal complexity is where the sum is minimal, which the train-test gap of gradient boosting against the number of stages makes visible.

When the classes are weighted, with weight $w_c$ for class $c$, the minimized loss is $\sum_iw_{c_i}\ell_i$ and the classifier estimates $\tilde p=w_sp/(w_sp+w_b(1-p))$ instead of $p$. A monotonic recalibration (isotonic regression~\cite{zadrozny_elkan}) restores the probability scale without changing the ranking, hence the AUC.

\subsection{Classification Without Labels}
Let two samples be mixtures of the same signal and background densities $s(x)$ and $b(x)$, with signal fractions $f_1\neq f_2$. The ratio of their feature densities, $p_1(x)/p_2(x)=[f_1\,s+(1-f_1)\,b]/[f_2\,s+(1-f_2)\,b]$, is a monotonic function of $s(x)/b(x)$, so that the classifier that is optimal between the two mixtures is, up to a monotonic transformation, the one that is optimal between signal and background~\cite{cwola}. The proof requires that $s$ and $b$ be the same in the two samples, which is the condition discussed in Sec.~\ref{sec:ml_labels}.

\subsection{Clustering and Uncertainty Estimators}
$K$-means~\cite{lloyd_kmeans} alternates the assignment of each point to the nearest centroid and the update of the centroids to the cluster means; each step lowers the inertia $W$, so the algorithm converges to a local minimum that depends on the initialization (hence several random starts). The adjusted Rand index~\cite{hubert_ari} between two partitions is $(\mathrm{RI}-E[\mathrm{RI}])/(\max\mathrm{RI}-E[\mathrm{RI}])$: it is $1$ for identical partitions and $0$ on average for random ones, and it is used to quantify the stability of the partition under a change of the initialization. The spread of the predictions of an ensemble, $\sigma^2(x)=\frac1M\sum_m[p_m(x)-\bar p(x)]^2$, estimates the epistemic part of the uncertainty when the members differ in the training data (bootstrap, as in the Random Forest) or in the initialization (deep ensemble~\cite{deep_ensembles}), and MC dropout~\cite{gal_mcdropout} obtains $M$ such members from a single network by keeping the dropout active at inference.
